\documentclass[10pt]{article}

\usepackage[utf8]{inputenc}

\usepackage[T1]{fontenc}

\usepackage{amsmath}

\usepackage{amsfonts}

\usepackage{amssymb}

\usepackage[version=4]{mhchem}

\usepackage{stmaryrd}

\usepackage{graphicx}

\usepackage[export]{adjustbox}

\graphicspath{ {./images/} }



\begin{document}

Вычисления в Р. м. сводятся к составлению следующей таблицы:



\[

\begin{array}{rrrrrr}

T_{00} & T_{01} & T_{02} & \ldots & T_{0, n-1} & T_{0 n} \\

T_{10} & T_{11} & T_{12} & \ldots & T_{1, n-1} & \\

\ldots & \ldots & \ldots & \ldots & & \\

T_{n-2,0} & T_{n-2,1} & T_{n-2,2} & & & \\

T_{n-1,0} & T_{n-1,1} & & & & \\

T_{n 0}, & & & & &

\end{array}

\]



где первый столбец состоит из квадратурных сумм (3) формулы трапеций. Элементы $(l+2)$-го столбца получаются из элементов $(l+1)$-го столбца по формуле $T_{k, l+1}=\frac{2^{2 l+2} T_{k+1, l^{-T} k_{l}}}{2^{2 l+2}-1}, k=0,1,2, \ldots, n-l-1$.



При составлении таблицы главная часть вычислительного труда затрачивается на вычисление элементов первого столбца. Элементы следующих столбцов вычисляются чуть сложнее конечных разностей.



Каждый элемент $T_{k l}$ таблицы есть квадратурная сумма, приближающая интеграл



\[

I \cong T_{k l} \text {. }

\]



Узлами квадратурной суммы $T_{k l}$ являются точки $j 2^{-k-l}, j=0,1,2, \ldots, 2^{k+l}$, а ее коэффициенты положительные числа. Квадратурная формула (5) точна для всех многочленов степени не выше $2 l+1$.



В предположении, что подинтегральная функция $f(x)$ имеет непрерывную производную порядка $2 l+2$ на $[0,1]$, разность $I-T_{k l}$ имеет представление вида (1), в к-ром $m=2 l+2$. Отсюда следует, что элементы $(l+2)$-го столбца, вычисляемые по формуле (4), являются улучшениями по Ричардсону элементов $(l+1)$-го столбца. В частности, для погрешности квадратурной формулы трапеций справедливо представление



\[

I-T_{k 0}=-\frac{f^{\prime \prime}(\xi)}{12} h^{2}, h=2-k, \xi \in[0,1],

\]



и способ Ричардсона дает более точное приближение к I:



\begin{center}

\includegraphics[max width=\textwidth]{2023_07_08_60f93b32c41b705d6edbg-1(1)}

\end{center}



$T_{k 1}$ оказывается квадратурной суммой формулы Симпсона, и т. к. Для погрешности этой формулы справедливо представление



\[

I-T_{k 1}=-\frac{f^{(\mathrm{I} \mathrm{V})}(\eta)}{180} h^{4}, h=2^{-k-1}, \eta \in[0,1],

\]



то снова можно воспользоваться способом Ричардсона и т. Д.



В Р. м. в качестве приближения к $I$ берется $T_{0 n}$, при этом предполагается, что существует непрерывная производная $f^{(2 n)}(x)$ на $[0,1]$. Ориентировочное представление о точности приближения $T_{0 n}$ можно получить, сравнивая $T_{0 n}$ и $T_{1, n-1}$.



Впервые метод изложен В. Ромбергом [1].



Jum.: [1] R o m be r g W., "Kgl. norske vid. selskabs forhandl.", 1955 , Bd 28 , № 7 , s. $30-36 ;$ [2] B a u e r F. L., R ut i s h a u s e r H., S t i e f e l E., "Proc. Symp. Appl. Math.»,

1963 , v. 15, p. $199-218$.

И. Mucosckux.



РОМБИЧЕСКАЯ СЕТЬ, к о н ф о р м о-ч е б ЫII е в с к а я с е т ь,- сеть, в каждом четырехугольнике к-рой, образованном двумя парами линий различных семейств, стороны равны с точностью до бесконечно малых 1-го порядка. На поверхности вращения ее асимптотическая сеть есть Р. с. Напр., на гиперболоиде вращения его прямолинейные образующие составляют P. c.



А. В. Иванов. РОСТА ИНДИКАТРИСА, и н д и к а т о р це ло й ф у в к ц и и, - величина



\[

h(\varphi)=\varlimsup_{r \rightarrow \infty} \frac{\ln \left|f\left(r e^{i \varphi}\right)\right|}{r^{\rho}},

\]



характеризующая рост целой функции $f(z)$ конечвого порядка $\rho>0$ и конечного типа $\sigma$ вдоль луча $\arg z=\varphi$ при больших $r\left(z=r e^{i \varphi}\right)$. Напр., для функции



\[

f(z)=e^{(a-i b)} z^{p}

\]



порядок равен $\rho$ и Р. и. равна $h(\varphi)=a \cos \rho \varphi+b \sin \rho \varphi$; для функции $\sin z$ порядок $\rho=1$ и $h(\varphi)=|\sin \varphi|$. Функция $h(\varphi)$ всюду конечна, непрерывна, имеет в каждой точке производную слева и справа, имеет производную всюду, кроме, быть может, счетного множества точек; всегда $h(\varphi) \leqslant \sigma$ и имеется, по крайней мере, одно $\varphi$, когда $h(\varphi)=\sigma$; обладает характерным с в о й с т в о. м т ри гономе т рич. вып уклости, т. ө. өсли



\[

h\left(\varphi_{1}\right) \leqslant H\left(\varphi_{1}\right), \quad h\left(\varphi_{2}\right) \leqslant H\left(\varphi_{2}\right)

\][



H(\textbackslash varphi)=a \textbackslash cos \textbackslash rho \textbackslash varphi+b \textbackslash sin \textbackslash rho \textbackslash varphi \textbackslash text \{, \}



][



\textbackslash varphi\_\{1\}<\textbackslash varphi\_\{2\}, \textbackslash varphi\_\{2\}-\textbackslash varphi\_\{1\}<\textbackslash min (\textbackslash pi / \textbackslash rho, 2 \textbackslash pi) \textbackslash text \{, \}

]



то $\quad h(\varphi) \leqslant H(\varphi), \varphi_{1} \leqslant \varphi \leqslant \varphi_{2}$.



Выполняется неравенство



\[

\left|f\left(r e^{i \varphi}\right)\right|<e^{[h(\varphi)+\varepsilon] r^{\rho}}, r>r_{0}(\varepsilon), \forall \varepsilon>0,

\]



где $r_{0}(\varepsilon)$ не зависит от $\varphi$.



P. и. вводится также для функции, аналитической в угле и имеющей в этом угле ковечный порядок вли уточненный порядок и конечный тип.



Лит.: [1] Л е вин Б. Я., Распределение корней целых функций, М., 1956 ; [2] M а р к у ш в в и ч А. И., Теория анаPOCTOK - термин, означающий точечную локализациюо различных математич. объектов (Р. функций, Р. отображений, Р. аналитических множеств и т. П.). Пусть, напр., $x$ есть точка топологич. пространства и $F$ - нек-рое семейство функций, определенных в окрестности $x$ (каждая в своей). Функции $f, g$ считаются эквивалентными (в точке $x$ ), если они совпадают в некрой окрестности $x$. Класс эквивалентностй по әтому отношению наз. р о с т ко м ф у н к ц и й класса $F$ в точке $x$. Так определяются Р. непрерывных функций, диффференцируемых функций в точках диффференцируемого многообразия, голоморфных функций в точках комплексного многообразия и т. п. Если семейство $F$ обладает нек-рой алгебраич. структурой, то множество Р. функций семейства $F$ наследует әту структуру (операции при помощи представителей классов). В частности, Р. голоморфных функций в точке $z$ образуют кольцо. Элементы поля частных этого кольца наз. р о с т -



\includegraphics[max width=\textwidth, center]{2023_07_08_60f93b32c41b705d6edbg-1}

Аналогично определяется Р. семейства подмножеств топологич. пространства. Напр., в точках аналитич. многообразия есть $P$. аналитических множеств (класс әквивалентности определяется по совпадению в окрестности данной точки). На Р . сөмейств подмножеств естественно определяются теоретико-множественные операции и отношения. Понятие Р. имеет смысл и для других объектов, определенных на открытых подмножествах топологич. пространства.



См. также Аналитическая функчия, Мероморфная бункция, Пучок.



Лum.: [1] Г а н н и н г Р., Р о с с и Х., Аналитические функции многих комплексных переменных, пер с англ., $\mathbf{M}$,

1969 .

E. чирка.



POTOP - то эке, что вихрь.



РУ ЛЕТТА - название плоской кривой, рассматриваемой как траєктория точки, жестко связанной с некрой кривой, катящейся по другой неподвижной кривой. В случае, когда окружность катится по прямой,





\end{document}